\documentclass[11pt,a4paper]{article}
\usepackage[top=2.0cm,bottom=2.0cm,left=2.2cm,right=2.2cm]{geometry}
\usepackage[T1]{fontenc}
\usepackage[utf8]{inputenc}
\usepackage{xcolor}
\usepackage{hyperref}

\pagestyle{empty}
\setlength{\parskip}{0.6em}
\setlength{\parindent}{0pt}

\definecolor{accent}{HTML}{1F4E78}
\hypersetup{colorlinks=true, urlcolor=accent, linkcolor=accent}

\begin{document}

\textbf{\large Prabhavit Gupta}\\
Via Zanetti 10, 30175 Venice, Italy ~$\vert$~ +39 351 610 7807 ~$\vert$~ \href{mailto:prabhavitg@gmail.com}{prabhavitg@gmail.com}\\
\href{https://www.linkedin.com/in/prabhavit/}{linkedin.com/in/prabhavit} ~$\vert$~ \href{https://github.com/prab-gupta}{github.com/prab-gupta}\\
26 September 2026

\textbf{Hiring Team}\\
Capgemini\\
Application for Data Engineer (Ref: 4374760503)

Dear Hiring Team,

I am writing to express my strong interest in the Data Engineer position at Capgemini. As an engineer specializing in data pipelines, modern data platform architectures, and scalable AI integrations based in Venice, I have followed Capgemini's pivotal role in global business and technology transformation. Capgemini's focus on designing Modern Data Platforms, implementing resilient Big Data architectures, and operationalizing scalable AI ecosystems for enterprise clients aligns directly with my practical background building automated ingestion pipelines, relational data architectures, and data validation systems.

Over the past three years across eight European companies, I have transformed complex operational requirements into reliable, production-ready data solutions. During my tenure at FlowSpace (London, remote), I led the core platform and data migration from a legacy architecture to a modern PostgreSQL, Next.js, and Laravel stack. In this role, I designed entity data models, implemented data mapping strategies, and executed automated migration scripts while scaling the platform from 2,000 to over 4,500 active users without service disruption.

To solve the challenge of ingesting heterogeneous information, I architected an automated data ingestion and ETL pipeline using TypeScript, n8n, and the Claude API. This system extracts, normalizes, and validates unstructured multi-format data (PDF brochures, spreadsheets, websites) into strictly typed backend schemas, with over 900 structured entities processed and human-in-the-loop quality controls. In addition, I engineered bidirectional REST integrations and webhook synchronization with OfficeRnD and Nexudus APIs to maintain data consistency across distributed booking records.

My daily engineering toolkit spans Python, SQL, TypeScript, Docker, and Git. I work comfortably with relational data modeling (PostgreSQL, MySQL), modern data warehousing principles, and asynchronous event-driven pipelines, alongside an understanding of distributed architectures such as Apache Spark and Kafka. I place a strong emphasis on data quality, automated verification, and clear stakeholder communication, translating complex technical architectures into tangible business outcomes for multidisciplinary teams.

Being based in Venice, I am in close proximity to Capgemini's Marcon office, enabling seamless hybrid collaboration in the Veneto area, with full availability for project travel to Milan, Rome, or client locations. I hold an Italian European Long-Term Residence Permit (EU citizen equivalent for work authorization, requiring no visa sponsorship) and communicate fluently in English (professional C2) with conversational Italian.

I would welcome the opportunity to discuss how my hands-on experience in automated data ingestion, relational schema engineering, and AI-enabled data pipelines can deliver immediate value to Capgemini and its enterprise clients.

Thank you for your time and consideration.

Sincerely,

Prabhavit Gupta

\vspace{0.6em}
{\footnotesize \textit{I hereby authorize the processing of my personal data contained in this document for the purposes of recruitment and selection in compliance with the European General Data Protection Regulation (EU Regulation 2016/679) and Italian Legislative Decree 196/2003, as modified by Legislative Decree 101/2018.}\par}

\end{document}
